\section{Lesson 15 --- 19/11/2026}

\subsection{Constrained optimization problems}

We now move our focus to constrained optimization problems. In particular, we will solve related unconstrained problems and, under suitable assumptions, we will show that their solutions converge to a solution of the given constrained problem.

A constrained problem is defined as follows:
\begin{equation}\label{eq:constr-P}
\begin{cases}
\min f(x)\\
g_i(x) \le 0 & \forall\, i=1,\dots,m\\
h_j(x) = 0 & \forall\, j=1,\dots,p.
\end{cases}
\tag{P}
\end{equation}
Let $X = \{x \in \R^n : g(x) \le 0,\ h(x) = 0\}$ be the feasible set of \eqref{eq:constr-P}.

The methods for solving \eqref{eq:constr-P} can be divided into two classes:
\begin{itemize}
    \item \g{Primal methods}: the method is applied directly to the given problem \eqref{eq:constr-P} (e.g., methods of changing the variables, descent direction methods such as the projected gradient method or the Frank-Wolfe method).
    \item \g{Dual methods}: methods that use the dual of \eqref{eq:constr-P} or related properties (e.g., gradient methods for solving the dual problem, penalty methods).
\end{itemize}

\subsubsection{Changing variables method}

As an example of a method of changing variables, we consider a problem with linear equality constraints only. We observe that a constrained problem with linear equality constraints
\[
\begin{cases}
\min f(x)\\
Ax = b,
\end{cases}
\]
where $A$ is a $p \times n$ matrix with $\operatorname{rank}(A) = p$, is equivalent to an unconstrained problem. Indeed, write $A = (A_B, A_N)$ with $\det(A_B) \neq 0$, where $A_B$ is a $p \times p$ matrix. Setting $x = \begin{pmatrix} x_B\\ x_N \end{pmatrix}$, then $Ax = b$ is equivalent to
\[
A_B x_B + A_N x_N = b \quad\Longrightarrow\quad x_B = A_B^{-1}(b - A_N x_N),
\]
thus, eliminating the variables $x_B$,
\[
\begin{cases}
\min f(x)\\
Ax = b
\end{cases}
\quad\text{is equivalent to}\quad
\begin{cases}
\min f\big(A_B^{-1}(b - A_N x_N),\, x_N\big)\\
x_N \in \R^{n-p}.
\end{cases}
\]

\begin{remark}
If $f$ is convex, then the previous unconstrained problem is an unconstrained convex problem.
\end{remark}

\begin{example}
Consider
\[
\begin{cases}
\min x_1^2 + x_2^2 + x_3^2\\
x_1 + x_3 = 1\\
x_1 + x_2 - x_3 = 2.
\end{cases}
\]
Since $x_1 = 1 - x_3$ and $x_2 = 2 - x_1 + x_3 = 1 + 2x_3$, the original constrained problem is equivalent to the following unconstrained problem:
\[
\begin{cases}
\min (1 - x_3)^2 + (1 + 2x_3)^2 + x_3^2 = 6x_3^2 + 2x_3 + 2\\
x_3 \in \R.
\end{cases}
\]
Therefore, the optimal solution is $x_3 = -1/6$, $x_1 = 7/6$, $x_2 = 2/3$.
\end{example}

\subsubsection{Penalty methods}

The change of variables works only for linear equality constraints. For problems with inequality constraints we follow a different idea: we remove the constraints and add to the objective a term that penalizes the points outside the feasible set. Consider a constrained optimization problem
\begin{equation}\label{eq:pen-P}
\begin{cases}
\min f(x)\\
g_i(x) \le 0 & \forall\, i=1,\dots,m.
\end{cases}
\tag{P}
\end{equation}
Let $X = \{x \in \R^n : g(x) \le 0\}$ be the feasible set of \eqref{eq:pen-P}.

\begin{definition}[Quadratic penalty function]
The \g{quadratic penalty function} is
\[
p(x) = \sum_{i=1}^{m} \big(\max\{0, g_i(x)\}\big)^2.
\]
For $\varepsilon > 0$, the \g{penalized problem} is the unconstrained problem
\begin{equation}\label{eq:pen-Peps}
\begin{cases}
\displaystyle \min f(x) + \frac{1}{\varepsilon}\, p(x) := p_\varepsilon(x)\\[1ex]
x \in \R^n.
\end{cases}
\tag{$P_\varepsilon$}
\end{equation}
\end{definition}

Note that
\[
p_\varepsilon(x)
\begin{cases}
= f(x) & \text{if } x \in X,\\
> f(x) & \text{if } x \notin X.
\end{cases}
\]

\begin{proposition}\label{prop:penalty}
\leavevmode
\begin{enumerate}
    \item If $f$ and $g_i$ are continuously differentiable, then $p_\varepsilon$ is continuously differentiable and
    \[
    \nabla p_\varepsilon(x) = \nabla f(x) + \frac{2}{\varepsilon}\sum_{i=1}^{m} \max\{0, g_i(x)\}\, \nabla g_i(x).
    \]
    \item If $f$ and $g_i$ are convex, then $p_\varepsilon$ is convex.
    \item Any \eqref{eq:pen-Peps} is a relaxation of \eqref{eq:pen-P}, i.e., $v(P_\varepsilon) \le v(P)$ for any $\varepsilon > 0$.
    \item If $x_\varepsilon^*$ solves \eqref{eq:pen-Peps} and $x_\varepsilon^* \in X$, then $x_\varepsilon^*$ is optimal also for \eqref{eq:pen-P}.
    \item If $0 < \varepsilon_2 < \varepsilon_1$, then $v(P_{\varepsilon_1}) \le v(P_{\varepsilon_2})$.
\end{enumerate}
\end{proposition}

Points 3 and 4 suggest to solve a sequence of penalized problems with smaller and smaller $\varepsilon$, until the solution found is feasible.

\begin{algorithm}[htbp]
\caption{Penalty method}
Set $\varepsilon_0 > 0$, $\tau \in (0,1)$, $k = 0$\;
Find an optimal solution $x^k$ of the penalized problem $(P_{\varepsilon_k})$\;
\While{$x^k \notin X$}{
    $\varepsilon_{k+1} = \tau\, \varepsilon_k$\;
    $k = k + 1$\;
    Find an optimal solution $x^k$ of the penalized problem $(P_{\varepsilon_k})$\;
}
\end{algorithm}

\begin{theorem}[Convergence]
\leavevmode
\begin{itemize}
    \item If $f$ is coercive, then the sequence $\{x^k\}$ is bounded and any of its cluster points is an optimal solution of \eqref{eq:pen-P}.
    \item If $\{x^k\}$ converges to $x^*$, then $x^*$ is an optimal solution of \eqref{eq:pen-P}.
    \item If $\{x^k\}$ converges to $x^*$ and the gradients of the active constraints at $x^*$ are linearly independent, then $x^*$ is an optimal solution of \eqref{eq:pen-P} and the sequence of vectors $\{\lambda^k\}$ defined as
    \[
    \lambda_i^k := \frac{2}{\varepsilon_k} \max\{0, g_i(x^k)\}, \qquad i=1,\dots,m,
    \]
    converges to a vector $\lambda^*$ of KKT multipliers associated to $x^*$.
\end{itemize}
\end{theorem}

\begin{remark}
By point 5 of \cref{prop:penalty}, the sequence of the optimal values $v(P_{\varepsilon_k})$ generated by the penalty method is nondecreasing. Moreover, if $x_{\varepsilon_k}^*$ solves $(P_{\varepsilon_k})$ and $x_{\varepsilon_k}^* \notin X$, $k = 1,2,\dots$, then, for $\varepsilon_{k+1} < \varepsilon_k$, we have $v(P_{\varepsilon_{k+1}}) > v(P_{\varepsilon_k})$, since $p_{\varepsilon_{k+1}}(x) > p_{\varepsilon_k}(x)$ for all $x \notin X$.
\end{remark}
